% Bloom levels: remember, understand, apply, analyze, evaluate, create.
\begin{problema}\label{prob:20ee012}
Define an $F$ random variable as a ratio of two independent normalized chi-square variables. State $E(F)$ when the denominator degrees of freedom are $d_2>2$.
\end{problema}

\begin{problema}\label{prob:2e8078c}
Explain why the reciprocal of a random variable with distribution $F_{5,8}$ has distribution $F_{8,5}$.
\end{problema}

\begin{problema}\label{prob:b47df2a}
Two production lines are compared for equality of variance. The first sample has $n_1=13$ and $S_1^2=45$; the second has $n_2=11$ and $S_2^2=20$. Test equality of variances at $\alpha=0.05$ using a two-sided procedure and the upper critical value $F_{12,10,0.025}=3.62$.
\end{problema}

\begin{problema}\label{prob:23d903e}
Prove that if $T\sim t_\nu$, then $T^2\sim F_{1,\nu}$.
\end{problema}

\begin{problema}\label{prob:e04a795}
Two groups have $n_1=15$, $S_1^2=50$, $n_2=12$, and $S_2^2=48$. Evaluate whether an equal-variance assumption is reasonable before choosing between a pooled two-sample $t$ procedure and Welch's procedure.
\end{problema}

\begin{problema}\label{prob:6277cc2}
Create an original example of a variance-comparison test using an $F$ statistic. Include the context, hypotheses, sample variances, statistic, critical value, and conclusion.
\end{problema}

\begin{solproblema}[prob:20ee012]
If $U\sim\chi^2_{d_1}$ and $V\sim\chi^2_{d_2}$ are independent, then
\[
	F=\frac{U/d_1}{V/d_2}\sim F_{d_1,d_2}.
\]
For $d_2>2$,
\[
	E(F)=\frac{d_2}{d_2-2}.
\]
\end{solproblema}

\begin{solproblema}[prob:2e8078c]
If
\[
	F=\frac{U/5}{V/8}\sim F_{5,8},
\]
where $U\sim\chi^2_5$ and $V\sim\chi^2_8$ are independent, then
\[
	\frac1F=\frac{V/8}{U/5}\sim F_{8,5}.
\]
The reciprocal swaps the numerator and denominator degrees of freedom.
\end{solproblema}

\begin{solproblema}[prob:b47df2a]
Use the larger sample variance in the numerator:
\[
	F=\frac{45}{20}=2.25.
\]
The degrees of freedom are $12$ and $10$. Since $2.25<3.62$, do not reject equality of variances at the $5\%$ level.
\end{solproblema}

\begin{solproblema}[prob:23d903e]
Write
\[
	T=\frac{Z}{\sqrt{V/\nu}},
\]
where $Z\sim N(0,1)$, $V\sim\chi^2_\nu$, and $Z$ and $V$ are independent. Then
\[
	T^2=\frac{Z^2}{V/\nu}=\frac{\chi^2_1/1}{\chi^2_\nu/\nu}\sim F_{1,\nu}.
\]
\end{solproblema}

\begin{solproblema}[prob:e04a795]
The variance ratio is
\[
	F=\frac{50}{48}=1.042.
\]
With $df_1=14$ and $df_2=11$, this is far below a typical upper $5\%$ critical value near $3.36$. There is no practical evidence of unequal variances. A pooled procedure is reasonable under the usual normality and independence assumptions, although Welch's method remains valid and is often a safe default.
\end{solproblema}

\begin{solproblema}[prob:6277cc2]
Example: two blood-pressure instruments are compared. Instrument A gives $n_1=10$ readings with $S_1^2=4.2$; instrument B gives $n_2=8$ readings with $S_2^2=1.5$. Test
\[
	H_0:\sigma_A^2=\sigma_B^2,\qquad H_a:\sigma_A^2\ne\sigma_B^2.
\]
The statistic is
\[
	F=\frac{4.2}{1.5}=2.8,
\]
with $df_1=9$ and $df_2=7$. Using an upper two-sided critical value about $4.82$, do not reject $H_0$.
\end{solproblema}
